\documentclass[11pt]{article}
\usepackage[margin=1in]{geometry}
\usepackage{amsmath,booktabs,hyperref,longtable}
\title{Autoregressive Next-Word Text Generation}
\author{Madhavan Komandur (Reg. 230953298)}
\date{}
\begin{document}
\maketitle
\section{Problem Statement}
Build and compare Simple RNN, Bidirectional RNN, LSTM, and GRU models for
autoregressive (next-word) text generation on \emph{The War of the Worlds}.
\section{Dataset}
Project Gutenberg eBook \#36 (\url{https://www.gutenberg.org/files/36/36-0.txt}).
After cleaning: \texttilda70k tokens, vocabulary $\approx 3.1$k, split 80/10/10, context length 20.
\section{Method}
Perplexity is defined as $PPL=\exp(\mathcal{L})$, where $\mathcal{L}$ is mean token cross-entropy.
Sampling uses $P(w)\propto \exp(z_w/T)$ with $T\in\{0.5,1.0,1.5\}$.
\section{Results}
Training metrics are saved in \texttt{results/metrics.json}. Representative values
will be inserted after the experiment completes:
\[
\begin{array}{lcccc}
\toprule
\text{Model} & \text{Params} & \text{Test PPL} & \text{Test Acc} & \text{Time(s)} \\
\midrule
\text{SimpleRNN} & - & - & - & - \\
\text{BiRNN} & - & - & - & - \\
\text{LSTM} & - & - & - & - \\
\text{GRU} & - & - & - & - \\
\bottomrule
\end{array}
\]
\section{Conclusion}
See \texttt{results/} and the Jupyter notebook \texttt{text\_generation.ipynb}.
\end{document}
PY